\chapter*{代数几何基础：评论与勘误}
\addcontentsline{toc}{chapter}{代数几何基础：评论与勘误}
\markboth{作者的评论与勘误}{}
\chapterintro{Fondements de la géométrie algébrique. Commentaires}{Bourbaki 讨论班，第14年，1961/62，补篇；1962年5月}
\origpage{3}{C-01}
\section*{说明}
1957至1962年间，我们在 Bourbaki 讨论班上作了八篇涉及代数几何基础的报告。除第一篇外，这些报告均在\term{概形}{scheme}{schéma}的框架内写成。所述的全部结果都将在 Jean Dieudonné 与 Alexander Grothendieck 合著的《代数几何原理》中得到安排。不过，这些报告中任何一篇的实质内容，目前都还没有被已出版或正在准备中的任何一章所涵盖，也没有被其他书籍或论文所涵盖；这种状况大概还会持续几年。这正是我们把这些报告汇编成册的主要原因：在正式的系统论述问世之前，先把概形理论中的若干关键概念和结果提供给读者。此外，阅读这些文字，可以迅速熟悉这些概念和结果，而不必一开始就陷入系统性专著中难免繁复的细节；同时也能获得研读这类专著所不可缺少的动机。

为方便读者，我们在此汇集若干评论和勘误，按各篇报告分组，特别指出自原文写成以来取得的进展，并补充一些参考文献。

这些报告中的若干结果，已在 IHÉS 的代数几何讨论班中得到详细讨论；1961/62年在 Harvard 先后举行的两个讨论班也作了详细讨论：第一个由我主持，第二个由 Mumford 与 Tate 主持。这两个讨论班的油印讲义由 Lichtenbaum 先生整理，目前正在准备中。

\subsection*{所用缩写}
\begin{description}
\item[SGA] \textit{Séminaire de géométrie algébrique}，Institut des Hautes Études Scientifiques，1960--1962（代数几何讨论班）。
\origpage{4}{C-02}
\item[SHG] \textit{Séminaire Grothendieck}，Harvard University，1961/62（Grothendieck 讨论班）。
\item[SHMT] \textit{Séminaire Mumford--Tate}，Harvard University，1961/62（Mumford--Tate 讨论班）。
\item[EGA] Grothendieck, A. 与 Dieudonné, J.，\textit{Éléments de géométrie algébrique}（《代数几何原理》），Paris，Presses universitaires de France，1960、1961、……；\textit{Institut des Hautes Études Scientifiques, Publications mathématiques}，4、8、11、……。
\item[TDTE] Grothendieck, Alexander，\textit{Technique de descente et théorèmes d'existence en géométrie algébrique}，I--VI（《下降技术与代数几何中的存在定理》），\textit{Séminaire Bourbaki}，第12卷，1959/60，第190、195号；第13卷，1960/61，第212、221号；第14卷，1961/62，第232、236号。
\end{description}

\section*{1. 凝聚代数层的对偶定理}
\noindent\textit{Séminaire Bourbaki}，第9卷，1956/57，第149号，25页。

\paragraph{关于第14页的评论。}正如我在1958年国际数学家大会的报告中所指出的，\footnote{Grothendieck, Alexander，\textit{The cohomology theory of abstract algebraic varieties}（抽象代数簇的上同调理论），\textit{Proceedings of the International Congress of Mathematicians [1958 Edinburgh]}，第103--118页，Cambridge，at the University Press，1960。}这里提出的问题现已完全解决。

关于\term{凝聚层}{coherent sheaf}{faisceau cohérent}的\term{对偶}{duality}{dualité}，读者还可以参看上述文献第112--115页、EGA III第二部分（准备中），以及在此之前可供参考的 SGA 1962。更系统的处理将放在 EGA 后续的一章中，按原计划是第IX章。

\section*{2. 形式几何与代数几何}
\noindent\textit{Séminaire Bourbaki}，第11卷，1958/59，第182号，28页。

\paragraph{第6页，定理6。}假设“$f$ 是\term{准射影}{quasi-projective}{quasi-projectif}的”可以替换为更弱的假设“$f$ 是\term{分离}{separated}{séparé}的”。这是由下述结果推出的（参看 SGA VIII，6.2）：每个\term{拟有限}{quasi-finite}{quasi-fini}且分离的\term{态射}{morphism}{morphisme} $f\colon X\to Y$ 都是准射影的。

\origpage{5}{C-03}
\paragraph{第12页及第14页的注1。}J.-P. Serre 已在\term{特征}{characteristic}{caractéristique} $p>0$ 的\term{代数闭域}{algebraically closed field}{corps algébriquement clos} $k$ 上，构造出一个三维的\term{非奇异射影簇}{nonsingular projective variety}{variété projective non singulière}，它不能由下述概形约化而来：该概形\term{固有}{proper}{propre}于一个\term{整局部环}{local integral domain}{anneau local intègre}，该环的\term{剩余域}{residue field}{corps résiduel}是 $k$，而其\term{分式域}{field of fractions}{corps de fractions}的特征为零。\footnote{Serre, Jean-Pierre，\textit{Exemples de variétés projectives en caractéristique $p$ non relevables en caractéristique zéro}（特征$p$下不能提升到特征零的射影簇的例子），\textit{Proc. Nat. Acad. Sc. U.S.A.}，47，1961，第108--109页。}据说 Mumford 也用非奇异射影曲面找到了类似的例子。

\paragraph{第15页，注2与注3。}对于这里猜测的结果，我现在已没有原先那么乐观。不过，注3末尾提出的关于\term{射影空间}{projective space}{espace projectif}结构变化的问题，Hironaka 已给予肯定解答；\term{阿贝尔簇}{abelian variety}{variété abélienne}的类似问题则由 Koizumi 解决。

\paragraph{第16页，第3行。}Mumford 最近已构造出\term{亏格}{genus}{genre}为 $g$ 的曲线的\term{模概形}{moduli scheme}{schéma des modules}（参看 SHMT）。此外，定理10证明，模理论中“$n$ 阶水平的 Jacobi 概形”是非奇异的，甚至在 $\mathbb Z$ 上\term{光滑}{smooth}{simple}。

\paragraph{第23页，第12行，即推论5的第一行。}补入：“$X$ 或 $Y$ 在 $k$ 上固有。”

第1至5节的实质内容已收入 EGA III 出版的部分，第6、7节的内容则收入 SGA III。关于\term{基本群}{fundamental group}{groupe fondamental}，参看 SGA V、IX、X、XI；关于 Lefschetz 型定理及许多未解决的问题，参看 SGA 1962 的第X、XII、XIII篇报告。只有\term{温和分歧覆盖}{tamely ramified covering}{revêtement modérément ramifié}的理论（参看定理14）尚未有正式成文的系统论述。定理14的推论完整地确定了代数闭域上代数曲线的、阶与特征互素的\term{伽罗瓦覆盖}{Galois covering}{revêtement galoisien}；它在下述三项工作中都起到了关键作用：
\begin{enumerate}[label=\arabic*°]
\item Igusa 对任意特征下非奇异射影曲面所满足的 Picard 不等式的证明；
\item Ogg 与 Šafarevič 各自独立开展的研究：考察定义在任意特征的一元\term{函数域}{function field}{corps de fonctions}上的阿贝尔簇，以及在该阿贝尔簇作用下的\term{主齐性纤维丛}{principal homogeneous bundle}{fibré principal homogène}所成的群；
\item Artin 最近对代数簇的“Weil 上同调”的若干关键定理的证明。
\end{enumerate}

\origpage{6}{C-04}
\section*{3. 下降技术与代数几何中的存在定理 I：一般性；忠实平坦态射的下降}
\noindent\textit{Séminaire Bourbaki}，第12卷，1959/60，第190号，29页。

\paragraph{第1页，第10行。}现在看来，说\term{下降}{descent}{descente}技术“是代数几何中大多数存在定理的基础”，未免言过其实。对于 TDTE I--VI 系列前两篇所讨论的非射影方法，这一说法在很大程度上确实成立；但对射影方法（TDTE IV、V、VI）则并非如此。

\paragraph{第5页，第16行。}不必假设 $\alpha$ 是一个 $\mathcal F$-下降态射。

\paragraph{第20页，注。}拟有限、\term{平展}{étale}{étale}且满射的态射 $S'\to S$，或者态射 $\Spec(\widehat A)\to\Spec(A)$，并不总是\term{严格下降态射}{strict descent morphism}{morphisme de descente strict}；即使 $A$ 是代数闭域 $k$ 上某条代数曲线的局部环，且 $S=\Spec(A)$，也仍然如此。例如，可以找到一个固有且光滑的态射 $f\colon X\to S$，使 $X$ 成为底为 $S$、在某条\term{椭圆曲线}{elliptic curve}{courbe elliptique} $E$ 作用下的主纤维丛，且
\[
f'\colon X\times_S S'\longrightarrow S'
\]
是射影的，而 $f$ 本身却不是射影的。因此，这同时也给出了阿贝尔概形作用下一个非\term{等平凡}{isotrivial}{isotrivial}的主齐性纤维丛。

\paragraph{第26页，第9行。}“Igusa--Lang”应为“Chow--Lang”。

下降理论的各种细节，参看 SGA VI、VII、VIII。

\section*{4. TDTE II：形式模理论中的存在定理}
\noindent\textit{Séminaire Bourbaki}，第12卷，1959/60，第195号，22页。

\paragraph{第8页。}命题5.1中的公式只有在 $\Lambda$ 是域时才正确。一般情形下，应把 $\mathfrak m_\xi/\mathfrak m_\xi^2$ 换成它对 $\mathfrak n_\xi/\mathfrak n_\xi^2$ 的像所取的商，其中 $\mathfrak n$ 是 $\Lambda$ 的\term{极大理想}{maximal ideal}{idéal maximal}。此外，文中关于 $O$ 在 $\Lambda$ 上光滑的定义，也只有在剩余域扩张 $k'/k$ 可分时才正确。一般情形参看 SGA III，1.1。

\paragraph{第14页，注。}这里提出的问题，在射影情形下已经由“\term{Hilbert概形}{Hilbert scheme}{schéma de Hilbert}”完全解决（参看 TDTE IV）。相反，Nagata 和 Hironaka 的例子表明：如果不作射影性假设，则所考虑的函子未必\term{可表}{representable}{représentable}；即便只考察某个三维非奇异完备簇的零维子簇的分类，也会如此。这与这样的簇的对称平方可能不存在有关。

\origpage{7}{C-05}
\paragraph{第15页，第3项。}更完整的讨论见 TDTE V。

\paragraph{第16页，命题3.1。}“如果 $f$ 固有”应为“如果 $f$ 固有且\term{可分}{separable}{séparable}”。

\paragraph{第18、19页，注的第2项。}我最近证明：域上阿贝尔簇的\term{形式模概形}{formal moduli scheme}{schéma formel des modules}确实在 Witt 环上光滑。换言之，设一个\term{阿廷局部环}{Artinian local ring}{anneau artinien local}是另一个阿廷局部环的商，则前者上的每个\term{阿贝尔概形}{abelian scheme}{schéma abélien}都可以由后者上的阿贝尔概形约化得到。证明只利用了第182号报告第12页所引入的\term{提升障碍类}{obstruction class to lifting}{classe d'obstruction au relèvement}的函子性。还应提及，亏格为 $g$ 的曲线及极化阿贝尔概形的模概形，已由 Mumford 构造出来（参看 SHMT）。

\paragraph{第19页，第5节。}用来定义 $U$ 的截面，不仅必须在 $\mathfrak X$ 上是\term{非零因子}{non-zero-divisor}{non diviseur de zéro}，还必须在每个 $X_n$ 上都是非零因子。

\paragraph{第20页，第12、13行。}“$G_x=G(\Spec(\widehat{\mathcal O}_{\mathfrak X,x'}))$”应改为
\[
G_x=G\bigl(\Spec(\mathcal O_{\mathfrak X,x'})\bigr)
\subset G\bigl(\Spec(\widehat{\mathcal O}_{\mathfrak X,x'})\bigr)；
\]
而 $\mathcal O'_{OX}$ 应改为 $\widehat{\mathcal O}'_{OX}$。

\section*{5. TDTE III：商预概形}
\noindent\textit{Séminaire Bourbaki}，第13卷，1960/61，第212号，20页。

\paragraph{第15页，断言(i)。}关于 $Y=X/R$ 在 $S$ 上准射影这一点，目前似乎只在 $\mathcal R$ 来自一个\term{等价关系}{equivalence relation}{relation d'équivalence}时才得到证明。

\paragraph{第18、19页。}正如我们在下一篇报告末尾所说，文中所提出的猜想的确是错误的。注8.1中提到的那个“肯定事实”，似乎由多位作者同时证明了，包括 Nagata、Rosenlicht、Grothendieck 等。

还应指出：TDTE V 利用本篇发展的取商理论来构造\term{Picard概形}{Picard scheme}{schéma de Picard}，这一应用也可以由对 Hilbert 概形的适当运用来代替（参看 SHMT）。正如第8节所说，这里所给理论的最大缺口，是缺少对非固有等价关系的商的存在判据，例如由射影群的某些作用产生的等价关系。
\origpage{8}{C-06}
Mumford 已在这一方向得到一个重要定理。\footnote{Mumford, David，\textit{An elementary theorem in geometric invariant theory}（几何不变量理论中的一个初等定理），\textit{Bull. Amer. Math. Soc.}，67，1961，第483--486页。}关于他的结果的改进及其在理论中的各种应用，参看 SHMT。

\section*{6. TDTE IV：Hilbert概形}
\noindent\textit{Séminaire Bourbaki}，第13卷，1960/61，第221号，28页。

\paragraph{第6页，定理2.1，第3行。}“必须”应改为“充分必要”。

\paragraph{第15页，第13行。}“$P\ne Q$”应改为“$P<Q$”。

\paragraph{同页，第15行。}“相反，如果 $P=Q$”应改为“其余情形下，对充分大的 $n$，有 $P(n)\ge Q(n)$”。

\paragraph{同页，第18行，公式$(*)$。}“$<$”和“$>$”应分别改为“$\le$”和“$\ge$”。

\paragraph{同页，倒数第5行。}相应符号应为 $Z$。

\paragraph{第18页，注3.9。}Hartshorne 已研究了代数闭域上 Hilbert 概形的\term{连通分支}{connected component}{composante connexe}。他证明 $\Hilb_r^P$ 连通，并确定了哪些数对 $(r,P)$ 满足 $\Hilb_r^P=\varnothing$。\footnote{Hartshorne, R.，\textit{Connectedness of the Hilbert scheme}（Hilbert概形的连通性），博士论文，Harvard，1962。}

\paragraph{第23页，倒数第5行。}“一个分支”应改为“一个分支的”。

\paragraph{第23、24页，注5.5。}关于 Zappa 的例子，还应提及：Mumford 刚刚构造了 $\mathbb P^3_{\mathbb C}$ 的 Hilbert 概形的一个\term{不可约分支}{irreducible component}{composante irréductible}，它在\term{泛点}{generic point}{point générique}处不是\term{约化}{reduced}{réduit}的；它的一般点表示次数14、亏格24的非奇异曲线。沿其中的一条曲线作\term{爆破}{blowing up}{éclatement}，他还得到一个 $\mathbb C$ 上三维的正则射影概形，其形式模概形在泛点处非约化。等价地说，按复解析几何意义理解的局部模簇，在每一点处都是非约化的（参看 \textit{Séminaire Cartan}，第13卷，1960/61）。

\origpage{9}{C-07}
\section*{7. TDTE V：Picard概形；存在定理}
\noindent\textit{Séminaire Bourbaki}，第14卷，1961/62，第232号，19页。

\paragraph{第7页，注3.3。}这里提出的问题，Mumford 已给予肯定解答（见对下一篇报告的评论）。

\paragraph{第13页，注5.2。}正如我在下一篇报告开头所说，这里提出的存在性猜想是错误的；不过，Mumford 用他的方法证明了一个稍弱的定理。

\paragraph{第14页，第7行。}在“完备”之后补入“且剩余域代数闭”。刚才提到的 Mumford 反例也表明，这一限制不可缺少。

\paragraph{第17页，注6.6。}正如我们在下一篇报告开头所说，这里提出的问题刚刚由 Murre 给予肯定解答。

\paragraph{第18页，注6.7。}这里提出的问题有肯定解答（见对下一篇报告的评论的最后一段）。

\paragraph{第18页，注6.8，第4行。}“\term{正则}{regular}{régulier}”应改为“在 $k$ 上光滑”。

\section*{8. TDTE VI：Picard概形；一般性质}
\noindent\textit{Séminaire Bourbaki}，第14卷，1961/62，第236号，23页。

\paragraph{第12页，第16行。}“\term{仿射直线}{affine line}{droite affine}”应改为“去掉原点的仿射直线”。

\paragraph{同页，第18行。}“$X[t]$”应改为“$X[t,t^{-1}]$”。

\paragraph{第16页，命题3.1，第6行。}“而假设……”应改为“而由假设……”。

\paragraph{第21页，第4节：Picard概形的有限性定理。}这一节提出的这类有限性问题，在本篇报告写成之后，差不多已全部解决。下面列出目前在这一方向上已知的主要事实。为使表述简洁，我们默认下述命题中出现的所有\term{Picard预概形}{Picard prescheme}{préschéma de Picard}均存在；其实只需对这些表述作显然的改写，就能完全去掉存在性假设。以下 $S$ 表示一个\term{诺特概形}{Noetherian scheme}{schéma noethérien}，$X$、$Y$ 表示 $S$ 上的固有概形。

\begin{statement}{(i)}
若 $f\colon X\to Y$ 为满态射，则
\[
f^*\colon \Pic_{Y/S}\longrightarrow\Pic_{X/S}
\]
是\term{有限型}{of finite type}{de type fini}态射。
\end{statement}

\origpage{10}{C-08}
\begin{statement}{(ii)}
设 $Y$ 在 $S$ 上射影，并配备一个相对于 $S$ \term{丰沛}{ample}{ample}的\term{可逆模}{invertible module}{Module inversible}；令 $X$ 是该模的任意一个截面的零点概形，$f\colon X\to Y$ 为典范浸入。再假设 $Y/S$ 的各纤维的不可约分支维数均至少为 $3$，则
\[
f^*\colon\Pic_{Y/S}\longrightarrow\Pic_{X/S}
\]
为有限型态射。
\end{statement}

\begin{statement}{(iii)}
设 $X$ 在 $S$ 上射影，且其所有\term{几何纤维}{geometric fibre}{fibre géométrique}均为维数 $n$ 的\term{整概形}{integral scheme}{schéma intègre}。设 $\mathcal O_X(1)$ 为 $X$ 上一个用来定义\term{Hilbert多项式}{Hilbert polynomial}{polynôme de Hilbert}的丰沛可逆模。若 $M$ 为 $\Pic_{X/S}$ 的一个子集，则 $M$ \term{拟紧}{quasi-compact}{quasi-compact}，当且仅当 $M$ 的元素的 Hilbert 多项式
\[
a_0x^n+a_1x^{n-1}+a_2x^{n-2}+\cdots
\]
中的系数 $a_1$ 和 $a_2$ 保持有界。
\end{statement}

\begin{statement}{(iv)}
对每个非零整数 $n$，群预概形 $\Pic_{X/S}$ 中的 $n$ 次幂同态都是有限型态射。
\end{statement}

注意，(i)、(ii)还意味着：在各自假设下，$\Pic_{Y/S}$ 的一个子集 $M$ 拟紧，当且仅当它在 $\Pic_{X/S}$ 中的像拟紧。因此，$Y$ 上的可逆模 $\mathcal L$ \term{$\tau$-等价于零}{$\tau$-equivalent to zero}{$\tau$-équivalent à zéro}，当且仅当它在 $X$ 上的逆像 $\tau$-等价于零。换言之，$\Pic^\tau_{Y/S}$ 是 $\Pic^\tau_{X/S}$ 的逆像。特别地，要证明前一个预概形在 $S$ 上有限型，只需对后一个证明这一点，因为 $\Pic_{Y/S}\to\Pic_{X/S}$ 为有限型。于是，利用(i)、Chow引理和(iii)，可得：

\begin{statement}{(v)}
$\Pic^\tau_{X/S}$ 在 $S$ 上有限型。
\end{statement}

一般而言，把有限态射的情形(i)与(ii)结合起来，就能把大部分有限性问题归结到 $X/S$ 的几何纤维为维数不超过 $2$ 的\term{整正规}{integral and normal}{intègre normal}概形的情形。通常，甚至还可以对不必有限的满态射使用(i)，再用代数曲面的\term{奇点解消}{resolution of singularities}{résolution des singularités}——任意特征下的结果由 Abhyankar 给出——把问题归结到 $X/S$ 光滑的情形，因此其几何纤维为维数 $2$ 的非奇异曲面。举例来说，结合(v)和用来控制非奇异射影曲面的 Néron--Severi 群之秩的 Picard--Igusa 不等式，这就能证明 Néron 有限性定理的下述推广：

\begin{statement}{(vi)}
设 $X/S$ 固有，除此之外不加限制。则 $X/S$ 的几何纤维 $X_i/k_i$ 的\term{Néron--Severi群}{Néron--Severi group}{groupe de Néron-Severi}
\[
\Pic_{X_i/k_i}/\Pic^0_{X_i/k_i}
\]
\origpage{11}{C-09}
均为有限生成群，且其秩以及挠子群的阶一致有上界。
\end{statement}

同样，归结到非奇异曲面的办法，再加上这一情形下已知的定理——即 Néron 有限性定理，以及 Néron--Severi 群上的\term{交形式}{intersection form}{forme intersection}非退化这一事实——还能推出：

\begin{statement}{(vii)}
设 $X$ 为代数闭域上的固有概形。则 $X$ 中存在有限条整闭曲线 $C_i$，$1\le i\le r$，具有下述性质：对 $\Pic_{X/k}$ 的任意子集 $M$，$M$ 拟紧，当且仅当整数
\[
\deg\bigl(\mathcal L|_{C_i'}\bigr),\qquad \mathcal L\in M,
\]
保持有界，其中 $C_i'$ 表示 $C_i$ 的\term{正规化}{normalization}{normalisée}。
\end{statement}

这里的 $r$ 可以取为 Néron--Severi 群的秩。一旦知道这个群有限生成，(vii)便归结为如下事实：由 $X$ 中各条曲线 $C$ 定义的、Néron--Severi 群上的线性形式，其共同零点恰是该群的挠元素。在 $X$ 射影且非奇异的情形，这一结果以及(v)均由 Matsusaka 给出。利用(vii)，容易得到 $X$ 上 $\tau$-等价于零的可逆模的下述刻画：

\begin{statement}{(viii)}
设 $X/k$ 为域上的固有概形，$\mathcal L$ 为 $X$ 上的可逆模。下述条件等价：
\begin{enumerate}[label=\alph*.]
\item $\mathcal L$ $\tau$-等价于零；
\item 对每个凝聚模 $F$，均有 $\chi(F\otimes\mathcal L)=\chi(F)$，其中 $\chi$ 表示\term{Euler--Poincaré示性数}{Euler--Poincaré characteristic}{caractéristique d'Euler-Poincaré}；
\item[\textnormal{b$'$.}] 条件同(b)，但只须取 $F=\mathcal O_Y$，其中 $Y$ 为 $X$ 中的一维整闭子概形；
\item[\textnormal{c.}] 对上述每个 $Y$，若以 $Y'$ 表示其正规化曲线，则有 $\deg(\mathcal L|_{Y'})=0$。
\end{enumerate}
若 $X/k$ 射影，且配备一个丰沛可逆模 $\mathcal O_X(1)$，则上述条件还等价于：
\begin{enumerate}[label=\alph*.,start=4]
\item 对每个整数 $m$，$\mathcal L^{\otimes m}(1)$ 都丰沛；
\item （$X$ 为整概形时。）对每一对整数 $m,n$，均有
\[
\chi\bigl(\mathcal L^{\otimes m}(n)\bigr)=\chi\bigl(\mathcal O(n)\bigr)。
\]
\end{enumerate}
\end{statement}
\origpage{12}{C-10}
也就是说，令 $F=\mathcal L^{\otimes m}(n)$ 时，条件(b)成立。

关于最后一个条件的充分性，应注意，它意味着各个 $\mathcal L^{\otimes m}$ 的 Hilbert 多项式全都相等。因此，由 Mumford 的判据(iii)，这些 $\mathcal L^{\otimes m}$ 落在 $\Pic_{X/k}$ 的一个拟紧子集中，亦即(a)成立。条件(b)、(b$'$)、(c)应视作\term{数值等价}{numerical equivalence}{équivalence numérique}概念在任意固有概形上的各种表述；通常，这个概念是对非奇异射影簇定义的。对后者，(a)与(c)的等价当然早已为人所知，归功于 Matsusaka。

判据(e)还推出下述结果：
\begin{statement}{(ix)}
设 $f\colon X\to S$ 为射影且\term{平坦}{flat}{plat}的态射，其几何纤维均为整概形。则 $\Pic^\tau_{X/S}$ 在 $\Pic_{X/S}$ 中既开又闭。
\end{statement}

对于关键结果(i)、(ii)、(iii)的证明，我们只作几句说明。结果(iv)与其他结果稍有不同；其证明只需对有限、满射且\term{根性}{radicial}{radiciel}的态射使用(i)，确切说，是对 Frobenius 态射使用(i)。证明(i)时，实质性地用到了非平坦下降的思想，特别参看 TDTE I 第9页。由于我们只寻求有限性结论，缺少\term{下降资料}{descent datum}{donnée de descente}的\term{有效性}{effectivity}{effectivité}判据并无妨碍。另一方面，Mumford 最近证明了(iii)的一个稍弱形式：其判据要求 Hilbert 多项式的所有系数有界。他的论证极其简单，受到 Nakai 的丰沛性判据之证明的启发；Nakai 原来只对非奇异曲面陈述了该判据，而 Mumford 已将它推广到任意射影态射。

不过，在我看来，这一论证还需要对 $X/S$ 的纤维附加一个稍强的限制，即 Serre 的 $(S_2)$ 性质；几何纤维正规时，这一性质成立。随后可把这一带限制的判据用于(ii)的证明：判据(i)允许我们归结到 $Y/S$ 平坦且几何纤维整、正规的情形，再使用 Mumford 的判据，便容易进一步使 $X/S$ 也满足相同条件。由维数假设可知，$Y$ 与 $X$ 的几何纤维在其闭点处的\term{深度}{depth}{profondeur}均至少为 $2$，因而可以使用 SGA 1962 第XII篇报告中所给形式的“等价判据”，从而完成(ii)的证明。有了(i)、(ii)以后，就不难在 Mumford 的判据中去掉对纤维正规性的任何假设，并把它证明到(iii)所述的更强形式。

\origpage{13}{C-11}
最后指出：(i)、(ii)的证明还表明，当 $S$ 是域 $k$ 的谱时，态射
\[
\Pic_{Y/k}\longrightarrow\Pic_{X/k}
\]
是\term{仿射}{affine}{affine}的，而不只是有限型的。
